\documentclass[UTF8]{ctexart}
\usepackage{xeCJK}
\usepackage{geometry}
\geometry{left=2cm,right=2cm,top=2.5cm,bottom=2.5cm}
\setlength{\headheight}{12.7pt}

\usepackage{titlesec}
\titleformat{\section}
  {\normalfont\large\bfseries}
  {\thesection}
  {1em}
  {}
\titlespacing*{\section}
  {0pt}
  {3.5ex plus 1ex minus .2ex}
  {2.3ex plus .2ex}

\usepackage{amsmath}
\usepackage{amssymb}
\usepackage{amsthm}
\usepackage{enumitem}
\setlist[enumerate]{itemsep=0pt,topsep=0pt,parsep=0pt,leftmargin=0pt,itemindent=2.5em}
\setlist[itemize]{itemsep=0pt,topsep=0pt,parsep=0pt,leftmargin=0pt,itemindent=2.5em}
\usepackage{fancyhdr}
\pagestyle{fancy}
\fancyhf{}
\usepackage{graphicx}
\usepackage{hyperref}
\usepackage{indentfirst}
\usepackage{lmodern}


\begin{document}
\setlength{\abovedisplayskip}{1pt}
\setlength{\belowdisplayskip}{1pt}

\section{无序求和}

\hypertarget{sec:定义1.1}{}
\noindent\textbf{定义1.1 (无序求和)}

给定半群$(A,+)$, 交换子集$B$, 非空有限集$X$和映射$f:X\to B$, 则存在唯一的$a\in A$使得
\vspace{3pt}
\[
    \forall\,\text{双射}\,\varphi:|X|\to X:\quad
    a=\sum_{i=0}^{|X|-1}f(\varphi(i)),
\]
把$a$记为$\displaystyle\sum_{x\in X}f(x)$, 有时也记作$\displaystyle\sum f$.

\noindent{\kaishu 验证.}

任取双射$\psi:|X|\to X$, 令$\displaystyle a=\sum_{i=0}^{|X|-1}f(\psi(i))$.
那么对任意的双射$\varphi:|X|\to X$,
\[
    \psi^{-1}\circ\varphi:|X|\to|X|
\]
是双射, 从而由\href{01 有序求和#nameddest=sec:定理2.2}{排列不变性定理}得
\vspace{3pt}
\begin{equation*}
    a=\sum_{i=0}^{|X|-1}f\psi(i)=\sum_{i=0}^{|X|-1}f\psi\psi^{-1}\varphi(i)
    =\sum_{i=0}^{|X|-1}f\varphi(i).
\tag*{\qedsymbol}\end{equation*}

\vspace{13pt}

特别地, 如果$(A,0_A,+)$是幺半群, 那么约定$\displaystyle\sum\varnothing=0_A$.

\vspace{13pt}

\hypertarget{sec:定理1.1}{}
\noindent\textbf{定理1.1}

任给半群$(A,+)$的交换子集$B$, 非空有限集$X,Y$, 双射$\sigma:Y\to X$
和映射$f:X\to B$, 有
\vspace{3pt}
\[
    \sum_{x\in X}f(x)=\sum_{y\in Y}f(\sigma(y)).
\]

\noindent{\kaishu 证明.}

取双射$\varphi:|X|\to X$, 则
\begin{equation*}
    \sum_{x\in X}f(x)=\sum_{i=0}^{|X|-1}f\varphi(i)=
    \sum_{i=0}^{|Y|-1}f\sigma\sigma^{-1}\varphi(i)=
    \sum_{y\in Y}f\sigma(y).
\tag*{\qedsymbol}\end{equation*}





\newpage

\hypertarget{sec:定理1.2}{}
\noindent\textbf{定理1.2}

任给半群$(A,+)$的交换子集$B$, 非空有限集$X$和映射$f:X\to B$.

\vspace{3pt}
如果非空的$X_1,X_2\subset X$且不相交, 那么
$\displaystyle\sum_{x\in X_1\cup X_2}f(x)=\sum_{x\in X_1}f(x)+\sum_{x\in X_2}f(x)$.

\noindent{\kaishu 证明.}

取双射$p:|X_1|\to X_1$和$q:|X_2|\to X_2$, 对任意的$i<|X|$定义
\[\varphi(i)=\left\{\begin{aligned}
    &p(i),&&i<|X_1|,\\[-3pt]
    &q(i-|X_1|),&&i\geqslant|X_1|,
\end{aligned}\right.\]
易证$\varphi:|X_1|+|X_2|\to X_1\cup X_2$是双射. 从而
\vspace{5pt}
\begin{equation*}\begin{aligned}
    \sum_{x\in X_1\cup X_2}f(x)&=\sum_{i=0}^{|X_1|+|X_2|-1}f\varphi(i)=
    \sum_{i=0}^{|X_1|-1}fp(i)+\sum_{i=|X_1|}^{|X_1|+|X_2|-1}fq(i-|X_1|)\\[3pt]
    &=\sum_{i=0}^{|X_1|-1}fp(i)+\sum_{i=0}^{|X_2|-1}fq(i)=
    \sum_{x\in X_1}f(x)+\sum_{x\in X_2}f(x).
\end{aligned}\tag*{\qedsymbol}\end{equation*}

\vspace{13pt}

\hypertarget{sec:定理1.3}{}
\noindent\textbf{定理1.3}

任给半群$(A,+)$的交换子集$B$, 非空有限集$X$和映射$f:X\to B$.

\vspace{3pt}
如果$S$是$X$的一个划分, 则$\displaystyle\sum_{x\in X}f(x)=
\sum_{s\in S}\sum_{x\in s}f(x)$.

\noindent{\kaishu 证明.}

取双射$p:|S|\to S$, 对任意的正整数$n<|S|$定义$\displaystyle X_n=\bigcup p[n]$,
显然有定义时
\vspace{1pt}
\begin{enumerate}
    \item $p[n+1]=p[n]\cup\{p(n)\}$, 且$p[n],\{p(n)\}$不相交,
    \vspace{1pt}
    \item $X_{n+1}=X_n\cup p(n)$, 且$X_n,p(n)$不相交,
\end{enumerate}
\vspace{3pt}
据此对正整数$n\leqslant|S|$归纳证明$\displaystyle
\sum_{x\in X_n}f(x)=\sum_{s\in p[n]}\sum_{x\in s}f(x)$.

首先, $n=1$时
\[
    \sum_{x\in X_1}f(x)=\sum_{x\in p(0)}f(x)=
    \sum_{s\in\{p(0)\}}\sum_{x\in s}f(x)=
    \sum_{s\in p[1]}\sum_{x\in s}f(x)\,;\\[3pt]
\]
其次, 假设命题对$n$成立且$n+1\leqslant|S|$, 那么
\vspace{3pt}
\[
    \sum_{x\in X_{n+1}}f(x)=\sum_{x\in X_n}f(x)+\sum_{x\in p(n)}f(x)=
    \sum_{s\in p[n]}\sum_{x\in s}f(x)+\sum_{s\in\{p(n)\}}\sum_{x\in s}f(x)=
    \sum_{s\in p[n+1]}\sum_{x\in s}f(x),\\[5pt]
\]
归纳成立. 取$n=|S|$即得$\displaystyle\sum_{x\in X}f(x)=
\sum_{s\in S}\sum_{x\in s}f(x)$. \hfill $\qedsymbol$

\end{document}